\section{Integrated Interpretation}

The four assignments fit as successive layers of the same preliminary missile design. The first layer fixes the vehicle scale and propulsion balance. The second checks whether the missile can be guided in the horizontal plane under the seeker limits. The third produces the aerodynamic and mass-property data needed for a 6DOF model. The fourth uses those data to size attitude and acceleration autopilots.

\subsection{Course Code Reuse}

\begin{itemize}
    \item Preliminary propulsion: the professor's sizing structure and \texttt{SimX} verification were retained.
    \item Point-mass guidance: the MANSUP Simulink model was reused, with edits focused on pre-programmed guidance and terminal switching.
    \item DATCOM: the processing chain with \texttt{DATCOMreader.m} and \texttt{Coef2Mat.m} was retained; the geometry input was rebuilt for the selected missile.
    \item 6DOF and autopilot: the course dynamic model and root-locus method were reused with the new missile data, CG/inertia and aerodynamic matrix.
\end{itemize}

\subsection{Main Engineering Decisions}

The decisions that most affected the final results were:

\begin{itemize}
    \item using solid propulsion only, requiring a long-burning solid sustainer;
    \item separating seeker activation from terminal-guidance authorization;
    \item using the same moment reference in DATCOM and 6DOF;
    \item correcting the DATCOM \(X_{CP}\) interpretation as a dimensionless value in calibers;
    \item tuning the acceleration autopilot gain around the root-locus result to satisfy both stability and bandwidth.
\end{itemize}

\subsection{Limitations}

The remaining gap is the integration between guidance and 6DOF control. The point-mass guidance model and the 6DOF/autopilot model were developed in separate course frameworks; a complete end-to-end Simulink mission remains the next integration step. Propulsion, structural masses and aerodynamic geometry are preliminary engineering estimates, and the autopilot was checked at representative points rather than over a full flight envelope. Within that scope, the data flow is consistent: geometry, mass properties, aerodynamic coefficients and control gains refer to the same missile configuration.

\subsection{Final Summary}

The final design package contains:

\begin{itemize}
    \item simulated range above \(70~\text{km}\);
    \item horizontal guidance logic for straight and off-boresight launches;
    \item DATCOM aerodynamic database for the equivalent geometry;
    \item section-based CG and inertia estimates;
    \item 6DOF dynamic checks;
    \item attitude and acceleration autopilots meeting \(5g\) and \(0.5~\text{Hz}\) requirements.
\end{itemize}
